\documentclass[11pt,letterpaper]{article}

\usepackage[margin=0.85in]{geometry}
\usepackage[T1]{fontenc}
\usepackage[utf8]{inputenc}
\usepackage[hidelinks]{hyperref}
\usepackage{array}
\usepackage{longtable}

\title{Keysight E36312A DC Sweep\\Manual de Usuario e Instalacion}
\author{ROMERUU-dev}
\date{2026-05-22}

\begin{document}
\maketitle
\tableofcontents
\newpage

\section{Objetivo}

Esta aplicacion controla una fuente Keysight E36312A/E36300 de tres canales por
PyVISA/SCPI. Permite control manual, barridos DC simples, barridos DC avanzados
con fuentes anidadas, graficas en tiempo real, exportacion de datos y apagado
seguro de salidas.

La version final de la interfaz esta orientada a instrumento real. El modo de
simulacion queda reservado para pruebas internas y no aparece en la interfaz de
usuario.

\section{Requisitos}

\subsection{Software}

\begin{itemize}
  \item Python 3.10 o superior.
  \item Git.
  \item Backend VISA: NI-VISA, Keysight IO Libraries o pyvisa-py.
  \item Dependencias Python listadas en \texttt{requirements.txt}.
\end{itemize}

\subsection{Hardware}

\begin{itemize}
  \item Fuente Keysight E36312A o modelo compatible E36300.
  \item Cable USB, LAN o GPIB segun la configuracion del laboratorio.
  \item Para mediciones de transistores: protoboard o fixture, cables cortos,
        resistencias de proteccion y un punto de GND comun.
\end{itemize}

\section{Instalacion desde GitHub}

En una PC nueva, clona el repositorio y crea un entorno virtual:

\begin{verbatim}
git clone https://github.com/ROMERUU-dev/keysight_E36312A_DCSweep.git
cd keysight_E36312A_DCSweep
python -m venv .venv
\end{verbatim}

En Linux o macOS:

\begin{verbatim}
source .venv/bin/activate
python -m pip install --upgrade pip
python -m pip install -r requirements.txt
\end{verbatim}

En Windows PowerShell:

\begin{verbatim}
.venv\Scripts\Activate.ps1
python -m pip install --upgrade pip
python -m pip install -r requirements.txt
\end{verbatim}

Si en Linux el comando \texttt{python} no existe antes de activar el entorno,
usa \texttt{python3}:

\begin{verbatim}
python3 -m venv .venv
source .venv/bin/activate
\end{verbatim}

\section{Ejecucion}

Modo recomendado:

\begin{verbatim}
python app.py
\end{verbatim}

La app abre maximizada, conservando la barra superior del sistema para mover,
cerrar, minimizar y maximizar. Atajos y opciones:

\begin{longtable}{>{\raggedright\arraybackslash}p{0.33\linewidth}p{0.57\linewidth}}
\textbf{Accion} & \textbf{Resultado} \\
\hline
\texttt{F11} & Entra o sale de pantalla completa real. \\
\texttt{Esc} & Sale de pantalla completa real. \\
\texttt{python app.py --windowed} & Abre en ventana normal sin maximizar. \\
\texttt{python app.py --fullscreen} & Abre directamente en pantalla completa real. \\
\texttt{python app.py --debug} & Activa logging mas detallado de SCPI y aplicacion. \\
\end{longtable}

Nota: en pantalla completa real, GNOME/Wayland oculta la barra superior de la
ventana. Para mover o cerrar con botones del sistema, sal de pantalla completa
con \texttt{F11} o \texttt{Esc}.

\section{Conexion al instrumento}

\begin{enumerate}
  \item Conecta la fuente Keysight por USB, LAN o GPIB.
  \item Abre la app.
  \item En la pestana \texttt{Connection}, presiona \texttt{Refresh}.
  \item Selecciona el recurso VISA detectado.
  \item Si no aparece, escribe el recurso manualmente, por ejemplo:
\end{enumerate}

\begin{verbatim}
USB0::0x2A8D::0x1102::<serial>::INSTR
\end{verbatim}

Despues presiona \texttt{Connect}. La app apaga las salidas al conectar y
muestra la respuesta de \texttt{*IDN?}. El ultimo recurso usado se guarda en
\texttt{config.json}, archivo local ignorado por Git.

\section{Permisos USBTMC en Linux}

Si Linux ve el dispositivo pero PyVISA no lista el recurso, revisa:

\begin{verbatim}
lsusb | grep -i keysight
ls -l /dev/usbtmc* /dev/bus/usb/*/* 2>/dev/null | grep -E 'usbtmc|2a8d' || true
\end{verbatim}

Regla udev sugerida:

\begin{verbatim}
sudo tee /etc/udev/rules.d/99-keysight-e36312a.rules >/dev/null <<'EOF'
SUBSYSTEM=="usb", ATTR{idVendor}=="2a8d", ATTR{idProduct}=="1102", GROUP="plugdev", MODE="0660"
KERNEL=="usbtmc*", ATTRS{idVendor}=="2a8d", ATTRS{idProduct}=="1102", GROUP="plugdev", MODE="0660"
EOF
sudo udevadm control --reload-rules
sudo udevadm trigger
\end{verbatim}

Desconecta y reconecta la fuente despues de aplicar la regla. Para una prueba
temporal, ajusta el bus/dispositivo real y usa:

\begin{verbatim}
sudo chmod g+rw /dev/usbtmc0 /dev/bus/usb/001/007
\end{verbatim}

\section{Pestana Connection}

\begin{itemize}
  \item \texttt{Refresh}: vuelve a buscar recursos VISA.
  \item \texttt{Connect}: abre el recurso seleccionado, apaga salidas y consulta
        \texttt{*IDN?}.
  \item \texttt{Disconnect}: ejecuta apagado seguro y cierra comunicacion.
  \item \texttt{Emergency Stop}: detiene barridos activos, intenta poner salidas
        en 0 V y apaga todos los canales.
\end{itemize}

\section{Pestana Manual}

La pestana \texttt{Manual} controla CH1, CH2 y CH3 individualmente.

\begin{itemize}
  \item \texttt{Voltage}: voltaje objetivo del canal.
  \item \texttt{Current limit}: limite de corriente.
  \item \texttt{Set V}: envia el voltaje.
  \item \texttt{Set I limit}: envia el limite de corriente.
  \item \texttt{ON/OFF}: prende o apaga la salida del canal.
  \item \texttt{Measure}: lee voltaje y corriente medidos.
  \item \texttt{All OFF}: apaga todos los canales.
  \item \texttt{Ramp to 0 V and OFF}: baja voltaje a 0 V y apaga salidas.
\end{itemize}

\section{Pestana Simple DC Sweep}

Esta pestana ejecuta un barrido de voltaje en un solo canal.

\begin{enumerate}
  \item Conecta el instrumento.
  \item Selecciona el canal.
  \item Define \texttt{V start}, \texttt{V stop}, \texttt{V step},
        \texttt{I limit}, \texttt{Settle time} y \texttt{Compliance tolerance}.
  \item Presiona \texttt{Start Sweep}.
  \item Observa la grafica en tiempo real.
  \item Presiona \texttt{Export CSV} si quieres guardar el barrido simple.
\end{enumerate}

El CSV de barrido simple se guarda por defecto en:

\begin{verbatim}
runs/sweeps/dc_sweep_latest.csv
\end{verbatim}

\section{Pestana Advanced DC Sweep}

\texttt{Advanced DC Sweep} permite configurar barridos tipo \texttt{.dc} con
hasta tres fuentes logicas. Cada fuente logica se asigna a un canal fisico.

\subsection{Orden de barrido}

\begin{itemize}
  \item \texttt{1st Source}: barrido rapido/interior.
  \item \texttt{2nd Source}: barrido exterior, util para familias de curvas.
  \item \texttt{3rd Source}: barrido mas externo.
\end{itemize}

Ejemplo:

\begin{verbatim}
.dc VDS 0 5 1 VGS 0 2 1
\end{verbatim}

La app ejecuta todos los valores de VDS con VGS=0, despues todos los valores de
VDS con VGS=1, y asi sucesivamente.

\subsection{Tipos de barrido}

\begin{itemize}
  \item \texttt{Linear}: inicio, fin e incremento lineal.
  \item \texttt{Decade}: puntos por decada.
  \item \texttt{Octave}: puntos por octava.
  \item \texttt{List}: lista manual de valores.
\end{itemize}

\subsection{Fuentes fijas}

La seccion \texttt{Fixed Sources} permite definir cada canal como:

\begin{itemize}
  \item \texttt{Off}: canal apagado.
  \item \texttt{Fixed voltage}: canal a voltaje fijo durante todo el barrido.
  \item \texttt{Swept source}: canal controlado por una fuente barrida.
\end{itemize}

La app valida que un canal no sea fuente fija y barrida al mismo tiempo.

\subsection{Presets}

\begin{itemize}
  \item \texttt{Single Channel I-V}: barrido I-V de un canal.
  \item \texttt{Diode I-V}: barrido bajo para diodos, con limite sugerido de
        5 mA a 10 mA.
  \item \texttt{NMOS Output Curves}: VDS en CH1 y VGS en CH2.
  \item \texttt{NMOS Transfer Curve}: VDS fijo en CH1 y VGS barrido en CH2.
  \item \texttt{BJT Output Curves}: VCE en CH1 y VBASE en CH2.
\end{itemize}

\section{Mediciones de transistores}

La pestana \texttt{Transistor} es una pantalla de etapa futura. Hoy las
mediciones reales de transistores se hacen desde \texttt{Advanced DC Sweep} con
presets.

\subsection{NMOS Output Curves}

Configuracion:

\begin{verbatim}
1st Source: VDS -> CH1, Linear, Start 0 V, Stop 5 V, Increment 0.05 V
2nd Source: VGS -> CH2, Linear, Start 0 V, Stop 3.3 V, Increment 0.1 V
3rd Source: disabled
X axis: Source1 value
Y axis: CH1 Imeas
Group curves by: Source2
\end{verbatim}

Conexion:

\begin{verbatim}
CH1+ -> Drain
CH1- -> GND comun
CH2+ -> Gate
CH2- -> GND comun
Source -> GND comun
\end{verbatim}

Interpretacion:

\begin{itemize}
  \item \texttt{ID} se aproxima con la corriente medida en CH1.
  \item \texttt{IG} se observa con la corriente medida en CH2.
  \item Cada valor de VGS produce una curva ID contra VDS.
\end{itemize}

\subsection{NMOS Transfer Curve}

Usa VDS fijo en CH1 y barre VGS en CH2. Grafica:

\begin{itemize}
  \item X: \texttt{Source1 value} o \texttt{CH2 Vmeas}.
  \item Y: \texttt{CH1 Imeas}.
  \item Group curves by: \texttt{None}.
\end{itemize}

\subsection{BJT Output Curves}

Usa CH1 como VCE y CH2 como VBASE. Debe existir una resistencia externa de base
para limitar corriente. La corriente de colector se observa como corriente en
CH1. No conectes la base directamente a una fuente sin resistencia de
proteccion.

\section{Archivos exportados}

Los barridos avanzados se exportan automaticamente en:

\begin{verbatim}
runs/YYYY-MM-DD_HH-MM-SS_<sweep_name>/
\end{verbatim}

Contenido:

\begin{longtable}{>{\raggedright\arraybackslash}p{0.30\linewidth}p{0.60\linewidth}}
\textbf{Archivo} & \textbf{Contenido} \\
\hline
\texttt{data.csv} & Datos por punto: setpoints, mediciones, potencia y compliance. \\
\texttt{metadata.json} & IDN, recurso VISA, tiempos, directiva .dc, fuentes y limites. \\
\texttt{sweep\_config.json} & Configuracion del barrido para repetirlo despues. \\
\texttt{plot.png} & Captura de la grafica. \\
\texttt{log.txt} & Eventos y condiciones de parada. \\
\end{longtable}

\section{Seguridad}

\begin{itemize}
  \item La app apaga salidas al conectar.
  \item Ante errores VISA/SCPI, intenta apagar salidas.
  \item Al cerrar, pregunta si debe poner salidas en 0 V y apagar.
  \item \texttt{Emergency Stop} debe usarse si el DUT entra en una condicion no deseada.
  \item Usa resistencias limitadoras y limites de corriente conservadores.
\end{itemize}

Limites por defecto:

\begin{longtable}{>{\raggedright\arraybackslash}p{0.20\linewidth}p{0.23\linewidth}p{0.23\linewidth}p{0.23\linewidth}}
\textbf{Canal} & \textbf{Voltaje} & \textbf{Corriente} & \textbf{Potencia} \\
\hline
CH1 & 0 a 6 V & 0 a 5 A & 30 W max \\
CH2 & 0 a 25 V & 0 a 1 A & 25 W max \\
CH3 & 0 a 25 V & 0 a 1 A & 25 W max \\
\end{longtable}

\section{Solucion de problemas}

\subsection{No aparece ningun recurso VISA}

\begin{itemize}
  \item Verifica cable y encendido del instrumento.
  \item En Linux, revisa permisos USBTMC y regla udev.
  \item Instala o revisa el backend VISA.
  \item Prueba escribir manualmente el recurso VISA.
\end{itemize}

\subsection{La ventana no muestra barra superior}

Eso ocurre solo en pantalla completa real. Presiona \texttt{F11} o \texttt{Esc}
para volver a una ventana normal/maximizada con barra del sistema.

\subsection{El barrido se detiene por compliance}

El DUT esta demandando corriente cercana al limite configurado. Reduce voltaje,
aumenta proteccion externa o ajusta el limite solo si es seguro para el circuito.

\subsection{No se genera plot.png}

El CSV y metadata son los archivos principales. Si la grafica no se guarda,
revisa permisos de escritura en \texttt{runs/}. La medicion no se invalida por
un fallo de captura de grafica.

\section{Desarrollo y pruebas}

Para correr pruebas:

\begin{verbatim}
source .venv/bin/activate
python -m pytest
\end{verbatim}

Las pruebas cubren generacion de rangos, validacion de seguridad, exportacion,
motor de barrido avanzado, layout principal y simulador interno.

\section{Actualizacion en otra PC}

Para actualizar una instalacion existente:

\begin{verbatim}
cd keysight_E36312A_DCSweep
git pull
source .venv/bin/activate
python -m pip install -r requirements.txt
python app.py
\end{verbatim}

\end{document}
